\section{Introduction}
\label{sec:introduction}

Population-based metaheuristics address problems for which exhaustive enumeration is impractical by combining diversification of the search space with intensification around promising solutions, and their behaviour depends on how these two activities are balanced. A useful balance need not remain constant during a run, since different stages of the search can benefit from different behaviours, and it can be regulated online from information generated by the search agents \cite{vega2021learning}. Machine learning has been brought into this regulation at the level of operators and parameters, with models built offline or updated online during the run \cite{talbi2021machine_tax}, and adaptive operator selection is among the integrations reviewed for combinatorial optimization \cite{karimimamaghan2022machine}.

Adaptive mechanisms differ in the signal used to drive that adjustment. A fixed schedule spends a predetermined budget on each phase without consulting the run, and a scalar stagnation counter, which measures only the time since the last improvement, cannot separate a temporary pause from a persistent plateau. Learned controllers represent richer relationships, but they require states, training and reward design before a policy can be applied \cite{adriaensen2022automated}. Selection hyperheuristics address a neighbouring question, choosing among low-level heuristics rather than adjusting a single algorithm \cite{drake2020recent}, and hyperheuristics in general operate on heuristics instead of searching the solution space directly \cite{dokeroglu2024hyper}. The search trajectory itself has also been exploited. Networks of visited solutions support the analysis and comparison of search behaviour \cite{ochoa2021search}, features computed from points already evaluated in a run can trigger a warm-started switch to another solver of a portfolio \cite{jankovic2022trajectory}, and trajectory features appear among the descriptors used for algorithm selection and configuration in black-box continuous optimization \cite{cenikj2026survey}.

These representations motivate a complementary question: how can the ordered best-so-far fitness curve provide feedback for controlling the joint parameter profile of a hybrid metaheuristic? This work proposes a training-free controller that uses the shape of the convergence curve as its control signal. A window of the most recent best-so-far values is compared against two synthetic references, a progress ramp and a stagnation plateau, using Dynamic Time Warping (DTW) \cite{berndt1994using} and its derivative variant DDTW \cite{keogh2001derivative}. Elastic alignment is relevant because improvements in a stochastic search do not occur at fixed instants, so two runs with the same qualitative pattern can have their jumps displaced in time. The controlled optimizer is a hybrid of the Grey Wolf Optimizer (GWO) and Differential Evolution (DE) operating on a single shared population, in which both operators generate a candidate for every individual and the better one survives. Rather than switching algorithms, the controller adjusts the joint parameter profile of the hybrid, so exploration and exploitation are regulated without discarding the population or the incumbent solution. Four interchangeable policies are defined over the same monitor: two alternate between fixed exploitation and exploration profiles, and two interpolate between them continuously as a function of the measured distances.

The controller is evaluated on the Multidimensional Knapsack Problem (MKP), using the OR-Library benchmark, where candidate solutions are mapped to binary vectors and repaired to restore feasibility. Beyond the discrete case, the same control layer is applied to a grid-connected wind--photovoltaic--electrolyzer--battery microgrid with green hydrogen production, a mixed continuous--discrete sizing problem built on an established techno-economic model. In both domains, the four adaptive policies are compared with two baselines that keep a fixed exploitation or exploration profile throughout the run. The comparison holds the operators, the binarization and the monitor settings constant, so that the contrast concerns the decision rule rather than a change of operator or monitor.

This study focuses on using the shape of the convergence curve as feedback for adjusting the parameter profile of a GWO--DE hybrid. The decision rule is training-free: it is not fitted to a data set and requires no training set, reward or offline policy search. Training-free is not parameter-free, however, since the window length and the percentile levels are chosen settings, and the thresholds adapt online only after a percentile history is available. This scope does not establish an advantage over learned, hyperheuristic or other adaptive parameter-control methods. The main contributions of this work are summarized as follows:

\begin{itemize}

\item A time-series view of convergence, in which the best-so-far fitness history is treated as an ordered temporal signal, together with a training-free decision rule that detects progress and stagnation by aligning this signal with reference curves through DTW or DDTW under run-relative percentile thresholds.

\item A shared-population GWO--DE hybrid whose joint $a$, $F$ and $CR$ profile is adjusted online from this signal, regulating exploration and exploitation without discarding the population or the incumbent solution.

\item A family of four control policies that share a single detector and differ only in how its output is used, either switching between fixed exploitation and exploration profiles or interpolating between them continuously.

\item A controlled evaluation on the MKP and on a green-hydrogen microgrid sizing problem, comparing the adaptive policies with fixed-profile baselines through non-parametric statistical tests and a one-at-a-time sensitivity analysis of the monitor settings.
\end{itemize}

The remainder of this paper is organized as follows. Section~\ref{sec:related_work} reviews hybrid coordination and the signals used to change search behaviour. Section~\ref{sec:background} introduces the knapsack formulation, the microgrid case, the GWO and DE operators, and the DTW monitor. Section~\ref{sec:proposed_approach} presents the shared-population hybrid and the four policies. Section~\ref{sec:experimental_setup} describes the experimental protocol. Section~\ref{sec:results} reports the campaigns, and Section~\ref{sec:statistical_analysis} compares the policies. Section~\ref{sec:discussion} interprets the findings. Section~\ref{sec:conclusions} presents the conclusions, limitations and future work.
